\subsection{Eksperimentell metode}
\label{sec:metode}
\subsubsection{Testformål, utstyr og underlag}
Testene skal dokumentere funksjon, kjøredistanse, repeterbarhet og retningsstabilitet. Underlag, måleutstyr og måleoppløsning beskrives før gjennomføring.
\subsubsection{Standardisert opptrekk og startprosedyre}
Trådlengde, opptrekksmåte, startposisjon, retning og frigjøring holdes så like som mulig. Lik trådlengde alene sikrer ikke lik tilført energi; variasjon i opptrekk registreres som en mulig feilkilde.
\subsubsection{Referansetester og sammenlignende forsøk}
Planen er å gjennomføre 10--20 referansekjøringer og minst 10 kjøringer per sammenlignet konfigurasjon. Endre én faktor om gangen der dette er praktisk, og registrer prototypeversjon og testrekkefølge. Definer måling av kjøredistanse og sideavvik entydig før testene.
\subsubsection{Registrering og behandling av data}
Loggen skal inneholde kjøringsnummer, dato, konfigurasjon, opptrekk, distanse, sideavvik, vellykket eller mislykket kjøring og observasjoner. Alle forsøk registreres; eventuelle utelatelser begrunnes. Rapporter gjennomsnitt, standardavvik, antall forsøk og suksessandel med tydelig angivelse av hvilke forsøk som inngår.
\subsubsection{Måleusikkerhet og feilkilder}
% Vurder opptrekk, underlag, startretning, måleavlesning og endringer i prototypen mellom kjøringene.
